% problem_id: S046-RD-cauchy-complete
% source_id: S046 (HoTT Book)
% source_file: reals.tex
% statement_locator: reals.tex lines 508--510
% proof_locator: reals.tex lines 512--538
% env: thm
% pairing: tight (verified)
% status: complete (statement + author proof verbatim + reason + locators)
%
\begin{thm} \label{RD-cauchy-complete}
  Every Cauchy approximation in $\RD$ has a limit.
\end{thm}

\begin{proof}
  Note that we are showing existence, not mere existence, of the limit.
  Given a Cauchy approximation $x : \Qp \to \RD$, define
  %
  \begin{align*}
    L_y(q) &\defeq \exis{\epsilon, \theta : \Qp} L_{x_\epsilon}(q + \epsilon + \theta),\\
    U_y(q) &\defeq \exis{\epsilon, \theta : \Qp} U_{x_\epsilon}(q - \epsilon - \theta).
  \end{align*}
  %
  It is clear that $L_y$ and $U_y$ are inhabited and rounded.  Disjointness follows from
  the Cauchy approximation condition. To establish locatedness, consider any 
  $q, r : \Q$ such that $q < r$. There is $\epsilon : \Qp$ such that 
  $4 \epsilon < r - q$. Since $q + 2 \epsilon < r - 2 \epsilon$ merely
  $L_{x_\epsilon}(q + 2 \epsilon)$ or $U_{x_\epsilon}(r - 2 \epsilon)$. In the first case
  we have $L_y(q)$ and in the second $U_y(r)$.

  To show that $y$ is the limit of $x$, consider any $\epsilon, \theta : \Qp$. Because
  $\Q$ is dense in $\RD$ there merely exist $q, r : \Q$ such that
  %
  \begin{narrowmultline*}
    x_\epsilon - \epsilon - \theta < q < x_\epsilon - \epsilon - \theta/2
    < x_\epsilon < \\
    x_\epsilon + \epsilon + \theta/2 < r < x_\epsilon + \epsilon + \theta,
  \end{narrowmultline*}
  %
  and thus $q < y < r$. It follows that $|y - x_\epsilon| < \epsilon + \theta$.
\end{proof}
